\documentclass[a4paper,11pt]{article}
\usepackage[T2A]{fontenc}
\usepackage[utf8]{inputenc}
\usepackage[russian]{babel}
\usepackage{amsmath,amssymb,mathtools}
\usepackage{icomma}
\usepackage[a4paper,top=22mm,bottom=22mm,left=24mm,right=24mm,headheight=22pt]{geometry}
\usepackage{microtype}
\usepackage{tikz}
\usetikzlibrary{arrows.meta,positioning,shapes.geometric,calc}
\usepackage{tabularx,booktabs}
\usepackage{enumitem}
\usepackage{fancyhdr}
\usepackage{xcolor}
\usepackage{hyperref}
\usepackage{needspace}
\definecolor{accent}{HTML}{008080}
\definecolor{darkaccent}{HTML}{005757}
\definecolor{textgray}{HTML}{353535}
\definecolor{softgray}{HTML}{737373}
\hypersetup{hidelinks,pdftitle={Чек-лист. Свойства логарифмов и вычисления},pdfauthor={Лёвин Артём Александрович}}
\setlength{\parindent}{0pt}
\setlength{\parskip}{5pt}
\setlength{\headsep}{9mm}
\raggedbottom
\setlist[enumerate]{leftmargin=*,itemsep=6pt,topsep=4pt}
\pagestyle{fancy}
\fancyhf{}
\fancyhead[L]{\small\color{textgray}Свойства логарифмов}
\fancyhead[R]{\small\color{textgray}Ярослав Гаврилов}
\fancyfoot[C]{\small\thepage}
\renewcommand{\headrulewidth}{0.55pt}
\renewcommand{\headrule}{\hbox to\headwidth{\color{accent}\leaders\hrule height \headrulewidth\hfill}}
\newcommand{\topic}[1]{\Needspace{5\baselineskip}\vspace{5pt}{\large\color{darkaccent}\bfseries #1}\par\vspace{4pt}}
\newcommand{\ex}[1]{\Needspace{4\baselineskip}\vspace{4pt}{\color{darkaccent}\bfseries Пример #1.}\enspace}
\newcommand{\key}[1]{\textcolor{darkaccent}{\textbf{#1}}}
\newcommand{\cond}{\ a>0,\ a\ne1,\ b>0}
\begin{document}
\thispagestyle{empty}
\begin{center}
\vspace*{25mm}
{\color{darkaccent}\large МАТЕМАТИКА}\par
\vspace{19mm}
{\LARGE Чек-лист}\par
\vspace{12mm}
{\color{accent}\Large\bfseries Свойства логарифмов\\и вычисление выражений}\par
\vspace{15mm}
{\large Подготовка к ЕГЭ по профильной математике}\par
\vspace{39mm}
{\large Ярослав Гаврилов}\par
\vspace{14mm}
{26 октября 2026 г.}\par
\vspace{16mm}
{\small Лёвин Артём Александрович\\[4pt]эксклюзивно для Ярослава Гаврилова}\par
\end{center}
\vfill
\begin{center}
\begin{tikzpicture}[x=1cm,y=1cm]
\draw[accent,line width=.7pt] (0,0) .. controls (1.1,.40) and (2.2,-.34) .. (3.5,0) .. controls (4.9,.35) and (5.8,-.33) .. (7.1,0);
\end{tikzpicture}
\end{center}
\clearpage

\topic{1. Что означает логарифм}
Для $a>0$, $a\ne1$, $b>0$ равенство
\[
\log_a b=x\quad\Longleftrightarrow\quad a^x=b
\]
означает: логарифм показывает, \emph{в какую степень нужно возвести основание $a$, чтобы получить $b$}.

\[\log_2 8=3\quad\text{поскольку}\quad 2^3=8,\qquad
  \log_a1=0,\qquad \log_a a=1.\]

\key{ОДЗ:} основание каждого логарифма положительно и не равно единице; его аргумент строго положителен. Эти условия проверяйте до преобразований.

\topic{2. Степень в аргументе логарифма}
\[
\boxed{\log_a(b^r)=r\log_a b},\qquad a>0,\ a\ne1,\ b>0.
\]
Показатель степени у \emph{аргумента} выходит вперёд без изменения. Формула работает в обе стороны: $r\log_a b=\log_a(b^r)$.

\ex{1}Вычислим логарифм, предварительно записав корень как степень:
\[
\begin{aligned}
\log_2\frac{1}{\sqrt[5]{8}}
 &=\log_2\left(2^{-3/5}\right)
 =-\frac35\log_2 2=-\frac35.
\end{aligned}
\]

\key{Полезный обратный ход:} любое действительное $x$ можно записать в виде
\[
 x=x\log_a a=\log_a(a^x),\qquad a>0,\ a\ne1.
\]
Например, $2=\log_3 9$, а $5=\log_{10}10^5$.

\topic{3. Степень в основании логарифма}
\[
\boxed{\log_{a^r} b=\frac1r\log_a b},\qquad
 a>0,\ a\ne1,\ b>0,\ r\ne0.
\]
Показатель степени у \emph{основания} выходит \emph{обратным множителем}. Если основание записано дробной степенью $a^{m/n}$, множитель равен $n/m$.

\ex{2}Здесь полезны сразу обе формулы:
\[
\begin{aligned}
\log_{\sqrt5}\frac{1}{\sqrt[3]{25}}
 &=\log_{5^{1/2}} 5^{-2/3}
 =\frac{1}{1/2}\cdot\left(-\frac23\right)\log_5 5
 =-\frac43.
\end{aligned}
\]
\clearpage

\topic{4. Сумма и разность логарифмов}
Для $a>0$, $a\ne1$, $b>0$, $c>0$:
\[
\boxed{\log_a b+\log_a c=\log_a(bc)},\qquad
\boxed{\log_a b-\log_a c=\log_a\frac bc}.
\]
Прежде чем объединять логарифмы, \key{приведите их к одному основанию}; если перед ними стоят коэффициенты, внесите их в показатели аргументов.

\ex{3}Объединим сумму в один логарифм:
\[
\log_2 3+2\log_2 5
=\log_2 3+\log_2 25
=\log_2 75.
\]

\ex{4}Основания различаются. Выразим восьмёрку степенью двойки:
\[
\begin{aligned}
\log_2 3+\log_8 5
 &=\log_2 3+\frac13\log_2 5\\
 &=\log_2 3+\log_2 5^{1/3}
 =\log_2\!\left(3\sqrt[3]{5}\right).
\end{aligned}
\]

\ex{5}Сочетание отрицательной степени основания и корня:
\[
\begin{aligned}
\log_{1/3}5+\log_{\sqrt[5]{27}}2
 &=\log_{3^{-1}}5+\log_{3^{3/5}}2\\
 &=-\log_3 5+\frac53\log_3 2\\
 &=\log_3\!\left(\frac{2^{5/3}}5\right).
\end{aligned}
\]

\topic{5. Основное логарифмическое тождество}
\[
\boxed{a^{\log_a b}=b},\qquad a>0,\ a\ne1,\ b>0.
\]
Для применения тождества основание степени и основание логарифма должны совпадать.

\ex{6}Если перед логарифмом стоит число, вносим его в аргумент:
\[
3^{2\log_3 5}=3^{\log_3 25}=25.
\]

\ex{7}Если в показателе есть сумма, используем свойства степеней:
\[
3^{1+\log_3 5}=3\cdot3^{\log_3 5}=3\cdot5=15.
\]
\clearpage

\topic{6. Если основания в степени и логарифме разные}
Используйте свойства степеней либо тождество
\[
\boxed{u^{\log_v w}=w^{\log_v u}},
\qquad u>0,\ w>0,\ v>0,\ v\ne1.
\]

\ex{8}Переходим к одному основанию:
\[
9^{\log_3 5}=(3^2)^{\log_3 5}
=3^{2\log_3 5}=5^2=25.
\]

\ex{9}В показателе разность, а основание логарифма равно $16=2^4$:
\[
\begin{aligned}
2^{3-\log_{16}5}
 &=\frac{2^3}{2^{\frac14\log_2 5}}
 =\frac{8}{\left(2^{\log_2 5}\right)^{1/4}}
 =\frac8{\sqrt[4]{5}}.
\end{aligned}
\]

\key{Обратное преобразование основания:} для положительных $p,q$ выполняется
\[
p=q^{\log_q p}\qquad(q\ne1).
\]
Это позволяет представить число степенью с выбранным основанием. Например, $3=5^{\log_5 3}$.

\topic{7. Переход к новому основанию}
Если $a,b,c>0$, $a\ne1$, $c\ne1$, то
\[
\boxed{\log_a b=\frac{\log_c b}{\log_c a}}.
\]
Частный случай при $b\ne1$ --- перемена мест аргумента и основания:
\[
\boxed{\log_a b=\frac1{\log_b a}}.
\]

\ex{10}Вычислим произведение, переписав единицы как логарифмы:
\[
\begin{aligned}
\left(\log_2 12-1\right)\left(\log_6 12-1\right)
 &=\log_2 6\cdot\log_6 2\\
 &=\log_2 6\cdot\frac1{\log_2 6}=1.
\end{aligned}
\]

\topic{8. С чего начинать вычисление}
\key{Главное правило занятия:} составные числа, корни и дроби сначала удобно представлять в степенном виде:
\[
27=3^3,\quad \sqrt[5]{125}=5^{3/5},\quad
\frac1{9}=3^{-2},\quad0{,}25=2^{-2}.
\]
Затем выберите нужное свойство логарифма, приведите основания и аккуратно выполните арифметику дробных показателей.
\clearpage

\begin{center}{\Large\color{darkaccent}\bfseries Алгоритм: преобразование выражений с логарифмами}\end{center}
\vspace{9mm}
\begin{center}
\begin{tikzpicture}[
    node distance=8mm,
    start/.style={draw=darkaccent,rounded corners=9pt,align=center,text width=8.1cm,minimum height=10mm,font=\small,inner sep=5pt},
    action/.style={draw=accent,rounded corners=3pt,align=center,text width=8.1cm,minimum height=12mm,font=\small,inner sep=5pt},
    decision/.style={draw=accent,diamond,aspect=3,align=center,text width=4.5cm,font=\small,inner sep=2pt},
    side/.style={draw=accent,rounded corners=3pt,align=center,text width=3.1cm,minimum height=18mm,font=\small,inner sep=5pt},
    flow/.style={-Latex,thick,color=darkaccent}
]
\node[start] (begin) {Дано выражение с логарифмами};
\node[action,below=of begin] (odz) {Проверьте ОДЗ: основания $>0$, $\ne1$; аргументы $>0$};
\node[action,below=of odz] (powers) {Представьте корни, дроби и составные числа в виде степеней};
\node[decision,below=11mm of powers] (agree) {Основания согласованы для выбранной формулы?};
\node[side,right=15mm of agree] (convert) {Приведите основания\ к общему виду или используйте смену основания};
\node[action,below=14mm of agree] (coeffs) {Если нужно, внесите коэффициенты в аргументы логарифмов};
\node[action,below=of coeffs] (formula) {Примените свойство: сумму, разность, степень или тождество};
\node[action,below=of formula] (check) {Проверьте знаки, дробные показатели и итоговый ответ};
\node[start,below=of check] (end) {Результат};
\draw[flow] (begin) -- (odz);
\draw[flow] (odz) -- (powers);
\draw[flow] (powers) -- (agree);
\draw[flow] (agree) -- node[left,font=\small] {да} (coeffs);
\draw[flow] (agree.east) -- node[above,font=\small] {нет} (convert.west);
\draw[flow] (convert.south) |- (coeffs.east);
\draw[flow] (coeffs) -- (formula);
\draw[flow] (formula) -- (check);
\draw[flow] (check) -- (end);
\end{tikzpicture}
\end{center}
\clearpage

\topic{9. Типичные ошибки и самопроверка}
\begin{enumerate}[label=\textcolor{darkaccent}{\arabic*.},itemsep=8pt]
\item \key{Путают показатель аргумента и основания.}
\[\log_a b^r=r\log_a b,\qquad
\log_{a^r} b=\frac1r\log_a b.\]
Вторая формула требует обратного множителя.

\item \key{Складывают аргументы.} Верно: $\log_a b+\log_a c=\log_a(bc)$; выражение $\log_a(b+c)$ в общем случае другое.

\item \key{Объединяют логарифмы с разными основаниями.} Например, $\log_2 3+\log_8 5$ нельзя сразу превратить в $\log_2 15$.

\item \key{Забывают о коэффициенте.} В выражении $2\log_3 5$ сначала получите $\log_3 25$, если хотите использовать формулу суммы или основное тождество.

\item \key{Теряют знак отрицательной степени.} Дробь $1/9=3^{-2}$, а число $0{,}25=2^{-2}$.

\item \key{Ошибаются при сокращении дробных показателей.} В вычислении
\[
\log_{\sqrt[5]{125}}\!\left(\sqrt[3]{25}\right)^4
=\log_{5^{3/5}} 5^{8/3}
=\frac{8/3}{3/5}=\frac{40}{9}
\]
деление на $3/5$ означает умножение на $5/3$.

\item \key{Оформляют решение слишком тесно.} Записывайте отдельно представление в виде степеней, перенос показателей и упрощение. Исправляя строку, зачёркивайте ошибку одной чертой, чтобы сохранялся ход рассуждений.
\end{enumerate}

\topic{Краткая памятка}
\begin{tabularx}{\linewidth}{@{}p{.39\linewidth}X@{}}
\toprule
\textbf{Признак в выражении} & \textbf{Первое полезное действие}\\
\midrule
Корень, дробь, составное число & Перейти к степени одного основания.\\
Степень у аргумента & Вынести показатель без изменения.\\
Степень у основания & Вынести обратный показатель.\\
Сумма или разность логарифмов & Согласовать основания и коэффициенты.\\
Степень с логарифмом & Проверить возможность применить $a^{\log_a b}=b$.\\
Основание неудобно & Использовать переход к новому основанию.\\
\bottomrule
\end{tabularx}
\clearpage

\begin{center}{\Large\color{darkaccent}\bfseries Тренировочный блок}\par\vspace{3pt}{\small Путь А · 10 заданий · по возрастанию сложности}\end{center}
\vspace{6mm}
\textbf{Базовый уровень}
\begin{enumerate}[label=\arabic*.,itemsep=12pt]
\item Вычислите: $\displaystyle \log_2 64$.
\item Вычислите: $\displaystyle \log_3\frac{1}{\sqrt[4]{81}}$.
\item Вычислите: $\displaystyle \log_{\sqrt[3]{5}}25$.
\end{enumerate}
\vspace{3mm}
\textbf{Средний уровень}
\begin{enumerate}[label=\arabic*.,start=4,itemsep=12pt]
\item Вычислите: $\displaystyle \log_{\sqrt[5]{125}}\left(\sqrt[3]{25}\right)^2$.
\item Сверните в один логарифм: $\displaystyle \log_2 3+2\log_2 5$.
\item Сверните в один логарифм: $\displaystyle \log_7 15-\log_7 3$.
\item Сверните в один логарифм: $\displaystyle \log_{1/3}4+\log_{\sqrt[3]{9}}5$.
\end{enumerate}
\vspace{3mm}
\textbf{Повышенный уровень}
\begin{enumerate}[label=\arabic*.,start=8,itemsep=13pt]
\item Вычислите: $\displaystyle 9^{\log_3 5}$.
\item Упростите: $\displaystyle 2^{3-\log_4 7}$.
\item Вычислите: $\displaystyle \left(\log_3 30-1\right)\left(\log_{10}30-1\right)$.
\end{enumerate}
\vfill
\textcolor{softgray}{\small Записывайте промежуточные преобразования; для дробей и корней сначала используйте степенную запись.}
\clearpage

\begin{center}{\Large\color{darkaccent}\bfseries Ответы}\end{center}
\vspace{10mm}
\renewcommand{\arraystretch}{1.6}
\begin{center}
\begin{tabularx}{.75\linewidth}{@{}cX@{}}
\toprule
\textbf{№} & \textbf{Ответ}\\
\midrule
1 & $6$\\
2 & $-1$\\
3 & $6$\\
4 & $\dfrac{20}{9}$\\
5 & $\log_2 75$\\
6 & $\log_7 5$\\
7 & $\displaystyle\log_3\!\left(\dfrac{5^{3/2}}4\right)$\\
8 & $25$\\
9 & $\displaystyle\dfrac8{\sqrt7}$\\
10 & $1$\\
\bottomrule
\end{tabularx}
\end{center}
\vspace{13mm}
{\color{darkaccent}\large\bfseries Итоговая памятка}\par\vspace{5pt}
Основной маршрут при вычислении логарифмических выражений:
\[
\text{ОДЗ}\ \longrightarrow\ \text{степенной вид}\ \longrightarrow\
\text{согласование оснований}\ \longrightarrow\ \text{формула}. 
\]
Для сложных комбинаций выполняйте действия последовательно: сначала вынесение показателей, затем объединение логарифмов или применение основного тождества, в конце --- арифметическая проверка.

\vfill
\begin{center}{\small\color{softgray}Лёвин Артём Александрович · эксклюзивно для Ярослава Гаврилова}\end{center}
\end{document}
